\documentclass[11pt]{article}

\usepackage[margin=1in]{geometry}
\usepackage{amsmath,amssymb,booktabs,microtype}
\usepackage[hidelinks]{hyperref}
\usepackage{xurl}

\newcommand{\vA}{A}
\newcommand{\cB}{B}
\newcommand{\KCL}{\mathcal{F}}
\newcommand{\Ephys}{\mathcal{E}_{\mathrm{phys}}}
\newcommand{\Eold}{\mathcal{E}_{\mathrm{old}}}

\title{Physical KCL Energy for Bidirectionally Amplified DRNs\\
\large Derivation, implementation correction, and experimental consequences}
\author{Internal technical note}
\date{26 August 2026}

\begin{document}
\maketitle

\begin{abstract}
This note explains a mismatch between the circuit equations and the interaction
energy formerly used for bidirectionally amplified deep resistive networks
(DRNs).  The corrected energy has the physical DRN equilibria as its stationary
points, up to a positive diagonal rescaling of the Kirchhoff-current-law (KCL)
vector field.  The former implementation instead placed an extra current-gain
factor on the receiving-node voltage inside every squared branch drop.  It was
therefore a coherent energy model, and its numerical residual could decrease to
zero, but its stationary points did not in general solve the stated physical
KCL equations.  The difference vanishes when the current gain is one; it affects
the legacy setting $(A,B)=(4,0.25)$ but not the baseline $(1,1)$ or the proposed
setting $(4,1)$.
\end{abstract}

\section{Circuit convention and KCL vector field}

Let $v_i^{(\ell)}$ denote the voltage of node $i$ in layer $\ell$.  A branch
with conductance $g_{ij}^{(p)}$ joins a sending node in layer $p$ to a receiving
node in layer $p+1$.  The voltage-controlled voltage source forwards
$A v_i^{(p)}$ from a hidden sending node, whereas the first branch is driven
directly by the clamped input.  Define
\begin{equation}
  \alpha_p =
  \begin{cases}
    1, & p=0,\\
    A, & p\geq 1,
  \end{cases}
  \qquad
  r_{ij}^{(p)} = \alpha_p v_i^{(p)}-v_j^{(p+1)}.
  \label{eq:branch-drop}
\end{equation}
Thus $g_{ij}^{(p)}r_{ij}^{(p)}$ is the current delivered to the receiving
node.  The current-controlled current source reflects the outgoing branch
current into the sending-node KCL equation with gain $B$.  With
$\KCL^{(\ell)}$ defined as net current delivered to the nodes of layer $\ell$,
one branch contributes
\begin{align}
  \KCL_{j,\mathrm{recv}}^{(p+1)}
    &= g_{ij}^{(p)}r_{ij}^{(p)},
    \label{eq:kcl-recv}\\
  \KCL_{i,\mathrm{send}}^{(p)}
    &= -B g_{ij}^{(p)}r_{ij}^{(p)}, \qquad p\geq1.
    \label{eq:kcl-send}
\end{align}
The complete vector field is obtained by summing the incoming and outgoing
branch contributions and the local one-port contribution at every free node.
The physical equilibrium is therefore characterized by
\begin{equation}
  \KCL^{(\ell)}\!\left(v^{(1)},\ldots,v^{(L)}\right)=0,
  \qquad \ell=1,\ldots,L,
  \label{eq:kcl-zero}
\end{equation}
subject to the perfect-diode constraints at the hidden nodes.  The output layer
is linear.

\section{Energy whose stationary points are the physical equilibria}

For the free layer $\ell\geq1$, introduce the positive diagonal layer factor
\begin{equation}
  d_\ell=\left(\frac{B}{A}\right)^{\ell-1}.
  \label{eq:diagonal-factor}
\end{equation}
The corrected energy contribution of a branch whose sending layer is $p$ is
\begin{equation}
  \Ephys_{ij,p}
  =\frac{1}{2}\left(\frac{B}{A}\right)^p
   g_{ij}^{(p)}\left(\alpha_p v_i^{(p)}-v_j^{(p+1)}\right)^2.
  \label{eq:physical-edge-energy}
\end{equation}
The factor $(B/A)^p$ is essential: it makes the asymmetric circuit field
integrable after a positive diagonal rescaling.  Differentiating
Eq.~\eqref{eq:physical-edge-energy} with respect to its receiving voltage gives
\begin{equation}
  \frac{\partial \Ephys_{ij,p}}{\partial v_j^{(p+1)}}
  =-\left(\frac{B}{A}\right)^p g_{ij}^{(p)}r_{ij}^{(p)}
  =-d_{p+1}\KCL_{j,\mathrm{recv}}^{(p+1)}.
  \label{eq:recv-gradient}
\end{equation}
For a free sending layer $p\geq1$,
\begin{align}
  \frac{\partial \Ephys_{ij,p}}{\partial v_i^{(p)}}
  &=\left(\frac{B}{A}\right)^p A g_{ij}^{(p)}r_{ij}^{(p)} \\
  &=d_p B g_{ij}^{(p)}r_{ij}^{(p)}
   =-d_p\KCL_{i,\mathrm{send}}^{(p)}.
  \label{eq:send-gradient}
\end{align}
After summing all branches and retaining the corresponding one-port
terms or perfect-diode constraints, the state gradient satisfies
\begin{equation}
  \nabla_{v^{(\ell)}}\Ephys=-d_\ell\KCL^{(\ell)},
  \qquad \ell=1,\ldots,L.
  \label{eq:gradient-kcl}
\end{equation}
Every $d_\ell$ is strictly positive for $A,B>0$.  Consequently,
\begin{equation}
  \nabla_v\Ephys=0
  \quad\Longleftrightarrow\quad
  \KCL(v)=0.
  \label{eq:zeros-equivalent}
\end{equation}
The energy gradient is not literally the negative KCL vector field when
$A\neq B$; it is the negative KCL field multiplied by a positive diagonal
metric.  This distinction changes the numerical scale by layer but not the
equilibrium set.

\section{The former energy and the extra factor}

The former dense and convolutional interactions evaluated
\begin{equation}
  \Eold_{ij,p}
  =\frac{1}{2}\left(\frac{B}{A}\right)^p
   g_{ij}^{(p)}\left(\alpha_p v_i^{(p)}-Bv_j^{(p+1)}\right)^2.
  \label{eq:old-edge-energy}
\end{equation}
Compared with Eq.~\eqref{eq:physical-edge-energy}, the receiving-node voltage
inside the square was multiplied by $B$.  This alters both the cross term and
the receiving-node quadratic term.  In particular,
\begin{equation}
  \frac{\partial \Eold_{ij,p}}{\partial v_j^{(p+1)}}
  =-\left(\frac{B}{A}\right)^p
    B g_{ij}^{(p)}\left(\alpha_p v_i^{(p)}-Bv_j^{(p+1)}\right),
  \label{eq:old-gradient}
\end{equation}
which is not a diagonal multiple of the physical receiving-node current in
Eq.~\eqref{eq:kcl-recv} unless $B=1$ (apart from degenerate cases).

This explains the selective scope of the correction:
\begin{center}
\begin{tabular}{lccc}
\toprule
Setting & $A$ & $B$ & Energy changed?\\
\midrule
Unamplified baseline & 1 & 1 & No\\
Proposed voltage amplification & 4 & 1 & No\\
Legacy reciprocal-current convention & 4 & 0.25 & Yes\\
\bottomrule
\end{tabular}
\end{center}

\section{Why the residual could still decrease}

There were two distinct consistency questions:
\begin{enumerate}
  \item Does the coordinate solver minimize the energy returned by the code?
  \item Does that energy generate the KCL equations of the stated circuit?
\end{enumerate}
They are not equivalent.

Before the 16 June 2026 change, the energy evaluator contained
Eq.~\eqref{eq:old-edge-energy}, while the coordinate-update coefficients still
encoded the raw physical KCL equations.  The coordinate solver could approach a
physical KCL root even though the autograd gradient of the reported old energy
did not approach zero.  Commit \texttt{1ba78c73} made the coordinate
coefficients agree with the derivatives of the old energy; commit
\texttt{08cfede3} carried that behavior into the amplification branch.  After
that change, the energy and its residual behaved normally because the solver
was consistently minimizing $\Eold$.

The observed decrease of the residual was therefore real, but it answered only
the first question.  It did not establish Eq.~\eqref{eq:gradient-kcl} for the
physical circuit.  The present correction restores both properties at once:
the coordinate coefficients agree with autograd, and the resulting gradient is
the diagonally scaled negative physical KCL field.

This also resolves the apparent contradiction with earlier successful BPTT and
ReLU-to-DRN experiments.  BPTT can train a differentiable recurrent solver even
when the solver represents the wrong equilibrium model.  In addition, older
zero-shot ReLU conversions either used the passive $A=B=1$ case or bypassed the
affected interaction evaluator, while the differential signed-resistive path
already used its own physical energy.  Successful accuracy is evidence that an
implemented model is trainable; it is not, by itself, a circuit-equation test.

\section{Implementation correction and regression tests}

Commit \texttt{599551c6} applies the correction to both
\texttt{DenseResistive} and \texttt{ConvResistive}.  It retains the depth factor
$(B/A)^p$ and the sending-node scale $\alpha_p$, but removes $B$ from the
receiving-node voltage in four coupled places:
\begin{itemize}
  \item evaluation of the interaction energy;
  \item the linear coefficient of each coordinate update;
  \item the receiving-node quadratic coefficient; and
  \item the conductance derivative used by energy-based learning.
\end{itemize}
The convolutional interaction receives the identical algebraic change after
unfolding its input patches.  Perfect-diode constraints and the output cost are
unchanged.

The regression tests check complementary invariants.  Per-interaction tests
compare coordinate coefficients with the autograd derivative of the evaluated
energy for dense and convolutional branches.  A two-edge dense-chain test
constructs the KCL currents explicitly and verifies Eq.~\eqref{eq:gradient-kcl}
for the affected $(A,B)=(4,0.25)$ setting.  The Conv1 compatibility study then
tests the end-to-end consequence using the frozen ordinary-MNIST protocol.

\section{Measured consequence and rerun policy}

On Conv1, unchanged-checkpoint replay and exact-learning-rate retraining show
that the impact depends on the optimizer:
\begin{center}
\begin{tabular}{llrr}
\toprule
Optimizer & Comparison & Old validation & Corrected validation\\
\midrule
SGD & best-checkpoint replay & 95.92\% & 93.84\%\\
SGD & exact-rate retrain, best & 95.92\% & 91.78\%\\
Adam & best-checkpoint replay & 96.44\% & 96.12\%\\
Adam & exact-rate retrain, best & 96.44\% & 96.50\%\\
\bottomrule
\end{tabular}
\end{center}
Using the predeclared $0.5$ percentage-point diagnostic threshold, the Adam
result is compatible at this resolution whereas the SGD result changes
materially.  This empirical compatibility does not undo the mathematical
correction; it motivates an Adam-first replacement program.

The active operational inventory therefore contains the 18 affected legacy
Adam arms: three clean wide BPTT arms, three clean centered-EqProp arms, three
noisy centered-EqProp arms, and nine fixed-initialization bounded BPTT arms.
The three clean-wide legacy SGD arms remain mathematically affected but are
deferred.  The immediate next diagnostic repeats only the exact frozen Adam
learning rates for Conv2 and Conv3, sequentially on one Jean Zay GPU at a time.
No official MNIST test split is read.  The detailed execution list and status
are maintained in
\path{docs/legacy_physical_kcl_rerun_inventory.md}.

\section{Conclusion}

The fixed energy is not an alternative modeling choice: it is the scalar
potential whose stationary points coincide with the zeros of the stated
amplified-DRN KCL vector field.  The former code could exhibit decreasing
energy and residual because it consistently minimized a different potential.
The corrected implementation preserves that numerical consistency while also
restoring the physical equilibrium equations.  Only legacy evidence with
$B\neq1$ needs replacement; the current compute priority is the Adam evidence.

\end{document}
